\documentclass[11pt]{article}

\usepackage[top=1in,bottom=1in,left=1.25in,right=1.25in]{geometry}
\usepackage{amsmath,amssymb,amsthm}
\usepackage{graphicx}
\usepackage{booktabs}
\usepackage{hyperref}
\usepackage{xcolor}
\usepackage{natbib}

\hypersetup{
  colorlinks=true,
  linkcolor=blue!60!black,
  citecolor=blue!60!black,
  urlcolor=blue!60!black,
}

\newtheorem{definition}{Definition}
\newtheorem{proposition}{Proposition}
\newtheorem{remark}{Remark}

\DeclareMathOperator*{\argmax}{arg\,max}

\title{\textbf{Multi-Agent Bandit Learning in Dark Liquidity Pools:\\
Collision Policies, Emergent Diversification,\\
and Adaptive Smart Order Routing}}

\author{
  M.~Eghe Aiani\\
  \texttt{m.egheaiani@student.vu.nl}
}

\date{March 2026}

\begin{document}

\maketitle

\begin{abstract}
Dark liquidity pools present a canonical but under-studied instance of the
multi-armed bandit problem: an institutional smart order router must
repeatedly allocate order flow across venues whose execution quality is
unobservable in advance and whose dynamics are non-stationary.
We introduce a multi-agent multi-armed bandit framework that models
competing institutions routing to the same set of dark pools, with three
explicitly-modelled \emph{collision mechanisms} — pro-rata fill splitting,
toxic-flow cancellation, and winner-takes-all latency priority — that
determine how rewards are shared when agents simultaneously select the same
venue.
Through three research questions, we show: (RQ1)~the collision mechanism
is the dominant driver of aggregate market welfare, with toxic-flow
cancellation destroying welfare proportional to the collision rate while
the other mechanisms are largely welfare-neutral; (RQ2)~under toxic-flow
cancellation, UCB-based agents spontaneously specialise on different pools
without any communication, a Nash-like separation emerging from pure
bandit feedback; and (RQ3)~adaptive strategies --- particularly
Sliding-Window UCB --- substantially outperform fixed-routing baselines in
a simulation calibrated to 222 weeks of real FINRA ATS transparency data
for Apple Inc. (AAPL), with the gap widening after each intraday venue
quality regime shift.
Our framework bridges multi-agent bandit theory, dark pool microstructure,
and real-world execution data, providing both a theoretical rationale for
adaptive smart order routing and empirical grounding through publicly
available regulatory reporting.
\end{abstract}

\tableofcontents
\newpage

% ─────────────────────────────────────────────────────────────────────────────
\section{Introduction}
\label{sec:intro}
% ─────────────────────────────────────────────────────────────────────────────

Dark liquidity pools --- privately operated, pre-trade opaque trading
venues --- account for approximately 35--40\% of US equity volume and a
growing share of European equity turnover \citep{foley2016dark}.
Institutional investors route a substantial fraction of their order flow
through dark pools to reduce market impact, but the choice of which pool
to access at any given time is a sequential decision problem under
uncertainty: execution quality varies across venues, shifts intraday, and
is not observable before submitting an order.

This is structurally a multi-armed bandit problem \citep{robbins1952}.
A smart order router (SOR) must learn which dark pools offer the best
combination of fill rate and price improvement, using only the feedback
received after each routing decision.
Unlike the single-agent setting studied by \citet{ganchev2010} and
\citet{bernasconi2022}, real markets feature \emph{multiple} competing
institutions whose order flow may arrive simultaneously at the same pool.
When this happens --- a \emph{collision} --- the outcomes depend on the
pool's matching logic: a crossing network may split the fill pro-rata,
an adversarial-selection-aware pool may cancel all fills, and a
latency-sensitive venue rewards only the fastest participant.
This collision structure has no counterpart in the single-agent bandit
literature and creates strategic dependencies between agents that are
central to real execution quality.

\paragraph{Contributions.} This paper makes the following contributions:
\begin{enumerate}
\item We formalise the dark-pool SOR problem as a \emph{multi-agent
  multi-armed bandit with collision externalities}, defining three
  collision policies (pro-rata, toxic-flow, winner-takes-all) and a suite
  of welfare and diversification metrics.

\item We run controlled experiments across seven learning strategies
  (UCB, Thompson Sampling, Sliding-Window UCB, KL-UCB/UCB-V,
  $\varepsilon$-greedy, greedy, random), three collision policies, varying
  agent counts, and three arm configurations, identifying the collision
  mechanism as the primary welfare driver.

\item We demonstrate emergent venue specialisation: under toxic-flow
  cancellation, agents spontaneously anchor on different pools through
  pure bandit feedback, with no explicit coordination.

\item We calibrate the simulation to 222 weeks of real FINRA ATS
  Transparency data (AAPL, five venues, 2021--2026) and show that
  Sliding-Window UCB outperforms static routing under regime shifts.
\end{enumerate}

\paragraph{Paper structure.} Section~\ref{sec:background} reviews
related work. Section~\ref{sec:framework} defines the formal model.
Section~\ref{sec:strategies} describes the strategies.
Sections~\ref{sec:rq1}--\ref{sec:rq3} present the three experiments.
Section~\ref{sec:discussion} discusses limitations and future work.
Section~\ref{sec:conclusion} concludes.

% ─────────────────────────────────────────────────────────────────────────────
\section{Background and Related Work}
\label{sec:background}
% ─────────────────────────────────────────────────────────────────────────────

\subsection{Multi-Armed Bandits}

The multi-armed bandit problem was formalised by \citet{robbins1952} as a
model of sequential experimentation under uncertainty.
\citet{lai1985} established the logarithmic lower bound on cumulative
regret and showed that any consistent policy must incur $\Omega(\log T)$
regret.
\citet{auer2002} introduced UCB1, a deterministic index policy that
achieves $O(\log T)$ regret with a simple, closed-form exploration bonus.
\citet{thompson1933} proposed a Bayesian alternative in which the agent
samples from the posterior distribution over arm means; modern analyses
\citep{bubeck2012} confirm near-optimal regret guarantees.
For non-stationary environments, \citet{garivier2011} introduced
Sliding-Window UCB (SW-UCB) and Discounted UCB (D-UCB), which restrict
learning to recent observations and achieve near-optimal regret under
abrupt change-point models.

\subsection{Multi-Player Bandits}

Multi-player bandits with collision costs were studied by
\citet{besson2018} who formalised the collision model, proposed RandTopM
and MCTopM, and provided finite-time regret guarantees.
\citet{boursier2019} introduced SIC-MMAB, achieving the information-theoretic
lower bound through deliberate collision signalling.
\citet{bistritz2018} proposed a fully decentralised algorithm achieving
$O(\log^2 T)$ regret without any communication.
Crucially, all prior work assumes a \emph{fixed} collision outcome: one
agent receives the reward, the rest receive zero.
Our framework generalises this to three distinct collision policies
(Section~\ref{sec:framework}), modelling the diverse matching logic of
real dark pool venues.

\subsection{Non-Stationary Bandits}

Real dark pool execution quality shifts intraday as liquidity providers
adjust their participation schedules, information events occur, and
MiFID~II volume caps bind \citep{buti2017}.
\citet{garivier2011} formalise this as a piece-wise stationary bandit
where arm means change at unknown change-points.
SW-UCB with window size $W$ achieves regret $O(\sqrt{T \Upsilon_T \log T})$
where $\Upsilon_T$ is the number of breakpoints --- close to the minimax
lower bound.
This motivates our use of SW-UCB in RQ3 as the primary non-stationary
baseline.

\subsection{Dark Pool Smart Order Routing}

\citet{ganchev2010} was the first to frame dark pool venue selection as a
bandit problem, studying a single agent routing across multiple pools and
demonstrating sublinear regret in a simulation.
\citet{bernasconi2022} extended this to combinatorial bandits, allowing
simultaneous submission to multiple pools subject to capacity constraints,
but retained the single-agent assumption.
On the market microstructure side, \citet{zhu2014} characterises the
adverse selection externalities that arise when informed traders use dark
pools, motivating our ``toxic-flow cancellation'' collision policy.
\citet{foley2016dark} empirically documents the relationship between dark
pool activity and price discovery, and \citet{buti2017} studies the
equilibrium determinants of dark pool volume, providing the stylised facts
that calibrate our arm parameters.

\paragraph{Literature gap.}
Table~\ref{tab:litgap} positions our work relative to the closest papers.
No prior work combines all four elements: multi-agent competition,
explicit collision modelling, dark pool context, and adaptive learning.

\begin{table}[h]
\centering
\caption{Position in the literature. \checkmark = property present;
  $\circ$ = partial; $\times$ = absent.}
\label{tab:litgap}
\begin{tabular}{lcccc}
\toprule
Paper & Multi-agent & Collision model & Dark pool & Adaptive learning \\
\midrule
Ganchev et al.\ (2010)    & $\times$ & $\times$ & \checkmark & \checkmark \\
Bernasconi et al.\ (2022) & $\times$ & $\times$ & \checkmark & \checkmark \\
Besson \& Kaufmann (2018) & \checkmark & $\circ$ & $\times$ & \checkmark \\
Boursier \& Perchet (2019)& \checkmark & \checkmark & $\times$ & \checkmark \\
Zhu (2014)                & \checkmark & \checkmark & \checkmark & $\times$ \\
\textbf{This work}        & \checkmark & \checkmark & \checkmark & \checkmark \\
\bottomrule
\end{tabular}
\end{table}

% ─────────────────────────────────────────────────────────────────────────────
\section{Framework}
\label{sec:framework}
% ─────────────────────────────────────────────────────────────────────────────

\subsection{Formal Model}

Let $\mathcal{A} = \{1, \ldots, K\}$ be a set of $K$ dark pool venues
(arms) and $\mathcal{N} = \{1, \ldots, N\}$ a set of $N$ agents
(competing institutional SORs).
At each timestep $t$, each agent $i$ selects a venue $a_i(t) \in \mathcal{A}$.
The \emph{collision set} for arm $a$ at time $t$ is
$\mathcal{C}_a(t) = \{i : a_i(t) = a\}$.

Each arm $a$ has a latent mean execution quality $\mu_a \in \mathbb{R}^+$
and noise $\sigma_a > 0$.
The \emph{raw reward} from arm $a$ at time $t$ is drawn independently:
\begin{equation}
  \tilde{r}_a(t) \sim \mathcal{N}(\mu_a,\ \sigma_a^2).
  \label{eq:raw_reward}
\end{equation}
The actual reward received by each agent in $\mathcal{C}_a(t)$ depends on
the pool's collision policy.

\begin{remark}[Gaussian reward approximation]
In practice, dark pool execution quality is not a single continuous scalar:
fills are binary events (order filled or not) with price improvement
conditional on filling.
The Gaussian model treats execution quality as a composite of fill probability
and price improvement, capturing mean and variance but abstracting away
the binary fill structure.
This simplification is standard in bandit-to-execution-quality mappings
\citep{ganchev2010} and does not qualitatively affect the results;
it means welfare quantities should be interpreted as expected execution
quality indices rather than as fill probabilities directly.
\end{remark}

\subsection{Collision Policies}

\begin{definition}[Collision Policy]
A collision policy $\pi : \mathbb{R} \times \mathbb{Z}^+ \to \mathbb{R}^N$
maps a raw reward $\tilde{r}$ and the number of colliding agents $n =
|\mathcal{C}_a(t)|$ to a vector of individual payoffs summing to at most
$\tilde{r}$.
\end{definition}

We model three policies:

\textbf{Pro-rata (Linear Share).}
Represents crossing networks and price-improvement venues that split the
available liquidity equally among all matched parties:
\begin{equation}
  \pi_{\text{linear}}(\tilde{r}, n) = \left(\frac{\tilde{r}}{n},\ \ldots,\ \frac{\tilde{r}}{n}\right).
\end{equation}

\textbf{Toxic-flow cancellation (Zero on Collision).}
Represents pools that detect crowded order flow as a signal of adverse
selection and cancel all fills \citep{zhu2014}:
\begin{equation}
  \pi_{\text{zero}}(\tilde{r}, n) = (0,\ \ldots,\ 0) \quad \text{for } n > 1.
\end{equation}
For $n = 1$ (no collision), both policies reduce to $\pi(\tilde{r}, 1) = \tilde{r}$.

\textbf{Winner-takes-all.}
Represents latency-sensitive venues where only the fastest participant fills:
\begin{equation}
  \pi_{\text{wta}}(\tilde{r}, n) = \tilde{r} \cdot \mathbf{e}_{i^*},
  \quad i^* \sim \mathrm{Uniform}(\mathcal{C}_a(t)),
\end{equation}
where $\mathbf{e}_{i^*}$ is the unit vector for the winning agent.
The winner is drawn uniformly at random among colliding agents.
This is a modelling simplification: in real price-time priority venues,
the fastest agent wins \emph{consistently}, creating a persistent latency
hierarchy between participants.
We use uniform random selection to avoid requiring latency rankings as
a model primitive; expected welfare is equivalent in both formulations,
though the uniform model distributes the winner's gain symmetrically across
agents rather than concentrating it in the consistently fastest one.
The limitation is discussed in Section~\ref{sec:discussion}.

\subsection{Metrics}

\textbf{Market welfare} is the total reward accrued by all agents per
timestep:
\begin{equation}
  W(t) = \sum_{i=1}^{N} r_i(t).
\end{equation}

\textbf{Collision rate} is the fraction of timesteps on which at least
one arm is chosen by two or more agents:
\begin{equation}
  \text{CR}(T) = \frac{1}{T}\sum_{t=1}^{T} \mathbf{1}\!\left[\exists\, a :
    |\mathcal{C}_a(t)| > 1\right].
\end{equation}

\textbf{Arm selection entropy} for agent $i$ over $T$ steps measures
diversification:
\begin{equation}
  H_i = -\sum_{a=1}^{K} \hat{p}_{i,a} \log \hat{p}_{i,a},
  \quad \hat{p}_{i,a} = \frac{1}{T}\sum_{t=1}^{T}\mathbf{1}[a_i(t) = a].
\end{equation}

\textbf{Cumulative regret} for a single agent against an oracle that
always selects the best arm:
\begin{equation}
  R(T) = \sum_{t=1}^{T} \left(\mu^* - r_i(t)\right),
  \quad \mu^* = \max_{a \in \mathcal{A}} \mu_a.
\end{equation}

\subsection{Arm Configurations}

Three arm configurations are used in RQ1/RQ2 experiments:

\begin{itemize}
\item \textbf{Wide gap}: means $\{1, 2, 3, 4, 5\}$, SDs $\{1.0, 1.2, 1.0, 1.5, 1.0\}$. One clearly dominant arm.
\item \textbf{Tight gap}: means $\{2, 2.5, 3, 3.5, 4\}$, uniform SDs $1.0$. Similar quality across arms.
\item \textbf{Uniform}: means $\{3, 3, 3, 3, 3\}$, SDs $1.0$. All arms identical.
\end{itemize}

% ─────────────────────────────────────────────────────────────────────────────
\section{Learning Strategies}
\label{sec:strategies}
% ─────────────────────────────────────────────────────────────────────────────

\subsection{Adaptive Strategies}

\textbf{UCB (UCB1)} \citep{auer2002}.
Selects the arm with the highest upper confidence bound:
\begin{equation}
  a^*(t) = \argmax_{a \in \mathcal{A}}\ \hat{\mu}_a(t) + \sqrt{\frac{2 \log t}{n_a(t)}},
\end{equation}
where $\hat{\mu}_a(t)$ is the sample mean and $n_a(t)$ the pull count up
to time $t$.
To break symmetry among identical agents, each agent's exploration order
is randomised at initialisation: without this, all agents pull the same
arm first and collide permanently in the initial exploration phase.

\textbf{Thompson Sampling (TS)} \citep{thompson1933}.
Maintains a Gaussian posterior over each arm's mean. With a
$\mathcal{N}(0, \tau_0^2)$ prior (we use $\tau_0^2 = 100$, uninformative
at the reward scale), the posterior after $n$ observations is
$\mathcal{N}(\hat{\mu}_a, \tau_n^2)$ where
$\tau_n^{-2} = n + \tau_0^{-2}$.
At each step, the agent samples $\theta_a \sim \mathcal{N}(\hat{\mu}_a, \tau_n^2)$
for each arm and selects $\argmax_a \theta_a$.

\textbf{Sliding-Window UCB (SW-UCB)} \citep{garivier2011}.
Restricts learning to the most recent $W$ timesteps:
\begin{equation}
  a^*(t) = \argmax_{a \in \mathcal{A}}\ \hat{\mu}_a^W(t) + \sqrt{\frac{\xi \log \min(t, W)}{n_a^W(t)}},
\end{equation}
where $\hat{\mu}_a^W$ and $n_a^W$ are the windowed mean and count.
We use $W = 200$ and $\xi = 1$.
This is the key strategy for non-stationary environments (RQ3), where arm
quality shifts between sessions.

\textbf{KL-UCB / UCB-V.}
Variance-adaptive UCB using empirical arm variance:
\begin{equation}
  a^*(t) = \argmax_{a \in \mathcal{A}}\ \hat{\mu}_a + \sqrt{\frac{2(\hat{\sigma}_a^2 + \sigma_0^2)\log t}{n_a(t)}},
\end{equation}
where $\hat{\sigma}_a^2$ is the sample variance and $\sigma_0^2 = 1$ is a
prior floor that prevents the bonus collapsing to zero before variance is
estimated. This tightens the bonus for low-variance arms while increasing
exploration for volatile ones.

\textbf{$\varepsilon$-Greedy.}
With probability $\varepsilon$ selects a uniformly random arm; otherwise
selects the empirically best arm. We test $\varepsilon \in \{0.05, 0.20\}$.
$\varepsilon = 0$ (greedy) is included as a lower bound.

\subsection{Static Baselines (RQ3)}

\textbf{Fixed routing.} Always sends orders to the same pre-assigned
venue. We test Fixed(best) (arm with highest initial quality), Fixed(mid),
and Fixed(worst) as bracketing benchmarks.

\textbf{Round-Robin.} Cycles sequentially through all venues, distributing
order flow equally without using execution quality feedback.

% ─────────────────────────────────────────────────────────────────────────────
\section{RQ1: Collision Policies and Market Welfare}
\label{sec:rq1}
% ─────────────────────────────────────────────────────────────────────────────

\subsection{Setup}

We run $N \in \{2, 3, 5, 8\}$ UCB agents for $T = 2{,}000$ steps on
the wide-gap arm configuration, averaged over 30 independent seeds.
All bar charts report 95\% confidence intervals computed as
$1.96 \times \hat{\sigma}/\sqrt{30}$ across seeds.
The primary metrics are aggregate market welfare $W(T)/T$ and collision
rate CR$(T)$.
We additionally vary arm configuration and strategy type to assess
robustness.

\subsection{Results}

\begin{figure}[h]
  \centering
  \includegraphics[width=\linewidth]{figures/rq1_base}
  \caption{RQ1-A: Aggregate welfare per step under three collision policies
    ($N=3$ UCB agents, wide-gap arms, 30 seeds, error bars = 95\% CI).}
  \label{fig:rq1_base}
\end{figure}

\begin{figure}[h]
  \centering
  \includegraphics[width=\linewidth]{figures/rq1_scaling}
  \caption{RQ1-B: Welfare vs.\ agent count under each collision policy.}
  \label{fig:rq1_scaling}
\end{figure}

\begin{figure}[h]
  \centering
  \includegraphics[width=\linewidth]{figures/rq1_armcfg}
  \caption{RQ1-D: Welfare across three arm configurations.}
  \label{fig:rq1_armcfg}
\end{figure}

\begin{figure}[h]
  \centering
  \includegraphics[width=\linewidth]{figures/rq1_strategy}
  \caption{RQ1-E: Strategy $\times$ policy welfare heatmap.}
  \label{fig:rq1_strategy}
\end{figure}

\paragraph{R1: Collision policy dominates welfare.}
Under pro-rata and winner-takes-all, aggregate welfare remains close to
the single-agent optimum regardless of collision rate.
Under toxic-flow cancellation, welfare decreases approximately linearly
with the collision rate: for 3 agents, mean welfare drops from
$\approx 4.9$ (linear share) to $\approx 3.1$ (zero on collision) at
collision rate $\sim 0.40$.

\paragraph{R2: Welfare degradation scales with agent count.}
Under zero on collision, adding agents initially raises the collision rate
and suppresses welfare. The relationship is near-linear from $N=2$ to
$N=7$, then plateaus as the collision rate saturates near 1 and all agents
learn to explore other arms.
Under linear share and winner-takes-all, aggregate welfare is flat with
agent count because the total reward is always distributed (not destroyed).

\paragraph{R3: Effect largest with heterogeneous arms.}
With uniform arms, all venues are equivalent and collision is costless
under pro-rata. With a wide quality gap, agents herd on the best arm and
collision costs under zero-on-collision are severe.

\paragraph{R4: Strategy matters less than policy.}
Replacing UCB with random agents barely changes aggregate welfare under
pro-rata — the collision policy determines the outcome, not the learning
algorithm. Under zero on collision, UCB recovers faster because agents
learn to disperse, but the initial welfare loss is determined by the
collision policy.

\paragraph{Summary.} These results directly address real practitioner
concerns: a pool that screens out crowded flow (zero on collision) may
actively harm aggregate institutional execution quality, even when
individual incentives to exploit it remain. This is consistent with the
adverse selection analysis of \citet{zhu2014}.

% ─────────────────────────────────────────────────────────────────────────────
\section{RQ2: Emergent Diversification Without Communication}
\label{sec:rq2}
% ─────────────────────────────────────────────────────────────────────────────

\subsection{Setup}

We run $N = 3$ UCB agents for $T = 3{,}000$ steps under each of the two
primary collision policies (zero on collision and linear share), averaged
over 30 independent seeds.
All figures report 95\% confidence intervals across seeds.
We compare early-phase (first 25\% of steps) versus late-phase (last 25\%)
collision rates, final arm-selection entropy, and per-agent late-phase
venue distributions.

\subsection{Mechanism}

Under zero on collision, the feedback mechanism drives agents toward the
Nash equilibrium of the collision game through reward degradation alone,
without any communication or strategic reasoning:
\begin{enumerate}
\item Each agent's initial exploration follows an independently randomised
  arm order (a deliberate design choice; without it, all agents explore
  in the same sequence and collide at every step, masking any
  differentiation signal entirely).
\item When agents collide on the best arm, both receive zero reward.
  Repeated zero rewards suppress that arm's UCB value estimate.
\item As agents' collision histories diverge due to their different
  exploration orders, one agent collides less on the best arm and
  accumulates a higher value estimate for it; the other's estimate
  remains suppressed.
\item The agent with the higher estimate anchors on the best arm and
  receives consistently high rewards, reinforcing its commitment.
  The other agent, facing a suppressed estimate, finds the second-best
  arm increasingly attractive. The asymmetry compounds.
\end{enumerate}

The result is convergence to the unique pure-strategy Nash equilibrium
of the underlying collision game: mutual separation, since any overlap
under zero on collision is strictly dominated (both agents receive zero
while each could unilaterally deviate to a non-overlapping arm and
receive positive reward).

\begin{remark}[Role of randomised exploration]
The convergence result depends on the randomised exploration design.
Without independent exploration orders, all agents would pull arms in
identical sequences, guaranteeing systematic collisions throughout the
initial phase and preventing divergence.
The result should therefore be stated precisely as: \emph{given
independent exploration paths, collision feedback alone is sufficient
to drive venue separation.}
This does not imply that UCB agents learn to coordinate strategically ---
they learn only individual arm values. The separation emerges because
zero on collision makes the best arm appear low-value to any agent that
herds onto it.
\end{remark}

Under linear share, collision reduces individual reward but never zeros
the arm's estimated value. Both agents continue to learn the arm is good,
and herding persists throughout the episode.

\subsection{Results}

\begin{figure}[h]
  \centering
  \includegraphics[width=\linewidth]{figures/rq2_entropy}
  \caption{RQ2-A: Rolling arm-selection entropy over time for each strategy
    ($N=3$ agents, zero-on-collision, window=150).}
  \label{fig:rq2_entropy}
\end{figure}

\begin{figure}[h]
  \centering
  \includegraphics[width=\linewidth]{figures/rq2_dist}
  \caption{RQ2-B: Per-agent arm-selection distributions (early vs.\ late phase).}
  \label{fig:rq2_dist}
\end{figure}

\begin{figure}[h]
  \centering
  \includegraphics[width=\linewidth]{figures/rq2_crowding}
  \caption{RQ2-C: Crowding effect — welfare vs.\ agent-to-arm ratio $N/K$.}
  \label{fig:rq2_crowding}
\end{figure}

\begin{figure}[h]
  \centering
  \includegraphics[width=\linewidth]{figures/rq2_collision}
  \caption{RQ2-D: Rolling collision rate over time per strategy.}
  \label{fig:rq2_collision}
\end{figure}

\begin{figure}[h]
  \centering
  \includegraphics[width=\linewidth]{figures/rq2_spec}
  \caption{RQ2-E: Per-agent specialisation — late-phase arm distributions
    for each strategy.}
  \label{fig:rq2_spec}
\end{figure}

\paragraph{R5: Zero on collision drives Nash equilibrium separation.}
Late-phase collision rate under zero on collision: $\leq 5\%$, down from
$\sim 60\%$ in the early phase (30-seed average with 95\% CI).
Under linear share, collision rate remains $\sim 40\%$ throughout.
Late-phase venue distributions show agents anchoring on \emph{different}
top arms under zero on collision.
This is not emergent coordination in a strategic sense: agents do not
communicate or model each other.
Rather, zero on collision makes the best arm appear low-value to any
agent that crowds it, causing the system to converge to the unique
pure-strategy Nash equilibrium of the collision game (mutual separation).

\paragraph{R6: Diversification pressure requires herding cost.}
Under linear share, EG agents herd persistently ($> 0.9$ probability mass
on the same top arm). UCB agents diversify under both policies due to
uncertainty-driven exploration, but the rate is substantially faster under
zero on collision where each collision actively suppresses the crowded
arm's value estimate.

\paragraph{R7: Crowding threshold effect.}
At crowding ratio $N/K > 0.6$, zero-on-collision welfare collapses sharply
and the reward-degradation pressure on the best arm intensifies.
Below this threshold, agents can spread across arms naturally.
This matches the intuition of MiFID~II's dark trading volume caps:
above a critical market-share threshold, the externality becomes severe
enough to require regulatory intervention --- or, in our model, strong
enough that purely self-interested adaptive agents will self-regulate.

\paragraph{Summary.}
The specialisation result under zero on collision is novel relative to
the multi-player bandit literature, which typically treats zero reward on
collision as a permanent failure mode rather than a convergence mechanism.
Our result shows that zero on collision functions as a \emph{reward
degradation signal} that eliminates herding as a dominated strategy,
driving the system to the Nash equilibrium venue separation without any
communication or explicit coordination mechanism.

% ─────────────────────────────────────────────────────────────────────────────
\section{RQ3: Calibrated Backtesting of Adaptive SOR}
\label{sec:rq3}
% ─────────────────────────────────────────────────────────────────────────────

\subsection{Empirical Calibration from FINRA ATS Data}

We calibrate the arm parameters using the FINRA ATS Transparency API,
which publishes weekly dark pool trading volume by venue and security.
We query 222 weeks of AAPL data (2021-12-06 to 2026-03-02), yielding
6{,}653 records across 37 active venues.

We select the top 5 venues by mean weekly share volume, applying a
$\geq 150$-week stability filter to exclude venues that entered or exited
the market during the sample period.
Table~\ref{tab:calibration} shows the resulting arm parameters.

The quality score is mean weekly share volume normalised so that UBS ATS
(the highest-volume venue) equals 5.0.
The SD is the actual week-to-week standard deviation of each venue's share
volume, scaled proportionally.

\textbf{Calibration caveat.}
The arm parameters are calibrated to \emph{trading volume}, not to
execution quality metrics (fill rate, price improvement, or execution
shortfall).
Volume and execution quality are positively correlated in large-cap stocks
--- \citet{foley2016dark} and \citet{buti2017} document this relationship
--- but they are not the same thing.
A high-volume venue may offer worse price improvement if it attracts
informed order flow, and a lower-volume venue may offer superior fill
quality for block trades.
SEC Rule~605 data provides direct fill rate and price improvement
statistics per venue, but is not available in machine-readable aggregate
form.
The RQ3 results should therefore be interpreted as establishing the
\emph{directional} finding --- adaptive $>$ static under non-stationarity
--- in a simulation grounded in real activity data, not as a
calibrated prediction of live strategy performance based on true
execution quality.

\begin{table}[h]
\centering
\caption{Empirically calibrated arm parameters from FINRA ATS data.
  AAPL, 222 weeks (2021-12-06 to 2026-03-02). Top 5 venues by mean
  weekly volume. CV = coefficient of variation (std/mean).}
\label{tab:calibration}
\begin{tabular}{clccccc}
\toprule
Pool & Venue & MPID & Mean & SD & CV & Weeks \\
\midrule
0 & JPM-X               & JPMX & 2.426 & 0.452 & 0.49 & 221 \\
1 & Level ATS           & EBXL & 2.839 & 0.706 & 0.53 & 239 \\
2 & Sigma X2            & SGMT & 3.196 & 0.844 & 0.50 & 239 \\
3 & Intelligent Cross   & INCR & 3.796 & 1.187 & 0.52 & 239 \\
4 & UBS ATS             & UBSA & 5.000 & 1.476 & 0.42 & 239 \\
\bottomrule
\end{tabular}
\end{table}

CV values of 0.42--0.53 across all venues confirm that dark pool volume is
genuinely stochastic: the best arm (UBS ATS, mean~5.0, SD~1.476)
substantially overlaps with the second-best (Intelligent Cross,
mean~3.796, SD~1.187), so identifying the best pool requires non-trivial
exploration.

\subsection{Regime-Shift Model}

To model intraday liquidity regime shifts, we partition the episode into
sessions of length $L = 300$ steps and permute arm quality scores randomly
between sessions, preserving the set of quality levels but reassigning
them across venues.
Agents retain their learned state across sessions --- they must detect and
adapt to the shift, not restart from scratch.
This models empirically documented phenomena: liquidity provider schedule
changes, information events, and MiFID~II volume cap binding
\citep{buti2017}.

\begin{figure}[h]
  \centering
  \includegraphics[width=\linewidth]{figures/rq3_static}
  \caption{RQ3-A: Cumulative reward of adaptive strategies vs.\ static
    baselines in a stationary environment (calibrated FINRA arms).}
  \label{fig:rq3_static}
\end{figure}

\begin{figure}[h]
  \centering
  \includegraphics[width=\linewidth]{figures/rq3_regime}
  \caption{RQ3-B: Cumulative reward under regime shifts
    (10 sessions $\times$ 300 steps). Vertical dashed lines mark session boundaries.}
  \label{fig:rq3_regime}
\end{figure}

\begin{figure}[h]
  \centering
  \includegraphics[width=\linewidth]{figures/rq3_comp}
  \caption{RQ3-C: Competitive advantage — per-agent cumulative reward in
    mixed adaptive vs.\ static multi-agent experiments.}
  \label{fig:rq3_comp}
\end{figure}

\subsection{RQ3-A: Stationary Environment}

We run all eight adaptive strategies plus four static baselines for
$T = 2{,}000$ steps with a single agent on the calibrated arms, averaging
over 10 seeds.

In a stationary environment, Fixed(best) routing achieves near-oracle
execution quality by definition --- it always selects UBS ATS (UBSA).
Adaptive strategies converge to near-oracle after an initial exploration
phase; their advantage over Fixed(best) is the \emph{ability to discover}
the best arm without prior knowledge.
Random routing achieves $\sim 66\%$ of oracle; Fixed(worst) achieves
$\sim 49\%$.
UCB, TS, SW-UCB, and KL-UCB all achieve $> 90\%$ of oracle quality after
convergence, with KL-UCB performing best under the heterogeneous variance
structure (Table~\ref{tab:calibration}: UBS ATS has $3.5\times$ the SD of
JPM-X, which amplifies the variance-adaptive bonus).

\subsection{RQ3-B: Non-Stationary (Regime Shifts)}

We run 10 sessions $\times$ 300 steps ($T = 3{,}000$ total) for
UCB, TS, SW-UCB, KL-UCB, Fixed(best), and Round-Robin.

\paragraph{R9: SW-UCB adapts fastest.}
With window $W = 200 < L = 300$, SW-UCB discards all pre-shift
observations within 200 steps of each regime shift and correctly
re-identifies the new best pool.
UCB accumulates stale history: its arm estimates are biased by pre-shift
data for the remainder of the session, and in some cases UCB never fully
recovers within a 300-step session.
Fixed(best) routing never re-evaluates: each regime shift permanently
degrades execution quality for that venue-specific strategy, so cumulative
regret grows piecewise linearly with each shift.

\paragraph{R10: Cumulative regret gap widens over regime shifts.}
After 10 sessions, SW-UCB accumulates $\sim 30\%$ less regret than UCB
and $\sim 60\%$ less regret than Fixed(best).
TS falls between UCB and SW-UCB: its Gaussian posterior updates with each
new observation, allowing some adaptation, but old data in the running
mean biases estimates more than SW-UCB's clean window.

\subsection{RQ3-C: Competitive Advantage}

We run 3-agent mixed experiments under pro-rata (linear share) with regime
shifts, comparing:
\begin{itemize}
\item \textbf{Setup A}: UCB vs Thompson Sampling vs Fixed(best)
\item \textbf{Setup B}: UCB vs Thompson Sampling vs Round-Robin
\end{itemize}

In both setups, adaptive agents accumulate higher cumulative reward than
their static counterpart by the end of the run, with the gap widening
after each regime shift as static agents fail to re-route.
In Setup A, the adaptive agents collectively exploit the new best venue
within $\sim 200$ steps of each shift; Fixed(best) continues routing to
what was the best venue before the shift and accrues lower quality fills.

\paragraph{Limitation.}
These results are from a calibrated simulation, not from live execution
data.
As noted in Section~\ref{sec:rq3}, arm parameters are calibrated to
\emph{trading volume}, not execution quality metrics.
The results establish the directional finding --- adaptive $>$ static
under non-stationarity --- grounded in real activity data, not a
calibrated prediction of live performance.

% ─────────────────────────────────────────────────────────────────────────────
\section{Discussion}
\label{sec:discussion}
% ─────────────────────────────────────────────────────────────────────────────

\subsection{Contributions Relative to Prior Work}

Our framework is the first to combine all four elements identified in
Table~\ref{tab:litgap}.
The collision mechanism is not a modelling convenience --- it captures a
domain-specific phenomenon (adverse selection in dark pools) that has
real welfare consequences and no counterpart in lit markets.
The emergent specialisation result of RQ2 is a novel finding relative to
the multi-player bandit literature, which does not study zero-on-collision
as a coordination signal.
The RQ3 calibration extends the bandit-to-dark-pool literature by grounding
the simulation in 222 weeks of real FINRA data, moving from plausible
parameters to empirically-derived ones.

\subsection{Limitations}

\textbf{Simulation gap.} The model abstracts away order-level mechanics
(partial fills, minimum quantity conditions, spread schedules), intraday
seasonality, and the strategic response of liquidity providers to routing
patterns. A reinforcement learning model with a richer state space would
be needed to capture these effects.

\textbf{Winner-takes-all simplification.}
The winner-takes-all policy selects the winner uniformly at random among
colliding agents, whereas real price-time priority venues assign the fill
to the consistently fastest participant.
The uniform model captures the correct expected welfare but misses the
persistent latency stratification of real WTA venues: in practice,
a single agent (the fastest) wins every collision, while the others
consistently receive zero.
This affects the learning dynamics --- a consistently-losing agent has
stronger incentive to diversify than a randomly-losing one --- and
the welfare distribution across agents.
An extension with explicit latency rankings would sharpen the WTA results.

\textbf{Volume as quality proxy.} Weekly share volume is a noisy proxy for
execution quality. High-volume venues are not always the best execution
venues: BIDS ATS, which was excluded from the top 5 due to its small-trade
count focus (mean trade size $\sim 396$ shares vs $\sim 60$--$90$ for other
venues), may offer superior block execution quality despite lower aggregate
volume.
SEC Rule~605 data provides direct fill rate and price improvement
statistics but is not available in machine-readable aggregate form.

\textbf{Single security.} Calibration uses AAPL, a Tier~1 NMS stock with
exceptional liquidity. Results may not generalise to smaller-cap stocks
where dark pool dynamics differ substantially \citep{foley2016dark}.

\textbf{Single-agent RQ3.} RQ3-A and RQ3-B use a single-agent setup to
isolate learning quality from collision interactions. Real markets feature
simultaneous routing by many institutions; the interaction of adaptive
learning and collision externalities in a non-stationary, multi-agent
setting is a natural extension.

\subsection{Future Work}

\textbf{Communication extension.} RQ2 establishes that diversification
emerges from pure bandit feedback. A natural next question is whether
allowing limited communication (e.g., noisy observation of competitors'
routing decisions, as required by MiFID~II post-trade venue reporting)
accelerates, replaces, or disrupts this emergent coordination.
Modelling this requires careful specification: observable actions with
noise is the most empirically grounded formulation; shared reward estimates
raise coordination legality concerns.

\textbf{Real execution data.} SEC Rule~605 data and institutional broker
execution quality reports (available under MiFID~II RTS~27/28) would
enable genuine backtesting beyond the calibrated simulation of RQ3.

\textbf{Contextual bandits.} Execution quality depends on observable
context: stock volatility, bid-ask spread, time of day, venue volume cap
proximity.
LinUCB or neural bandit architectures could incorporate these signals,
turning the problem into a contextual bandit and potentially further
improving adaptive routing.

% ─────────────────────────────────────────────────────────────────────────────
\section{Conclusion}
\label{sec:conclusion}
% ─────────────────────────────────────────────────────────────────────────────

We have introduced a multi-agent multi-armed bandit framework for dark
liquidity pool trading that bridges bandit theory, market microstructure,
and real regulatory data.
Three main findings stand out.

First, the collision mechanism is the primary determinant of aggregate
market welfare: toxic-flow cancellation destroys welfare proportional to
the collision rate, while pro-rata and winner-takes-all policies are
largely welfare-neutral. This has direct implications for institutional
execution strategy: routing to pools that screen out crowded flow may
individually identify adverse selection but collectively harm market
welfare.

Second, under toxic-flow cancellation, competing adaptive agents
spontaneously specialise on different pools through pure bandit feedback
alone. This emergent coordination --- a Nash-like venue separation without
communication --- suggests that collision-aware pool design may
inadvertently foster the market structure equivalent of a Nash
equilibrium, potentially beneficial from a stability perspective.

Third, Sliding-Window UCB substantially outperforms both UCB and static
routing in a simulation calibrated to 222 weeks of real FINRA ATS data
for AAPL, with the advantage growing with each intraday venue quality
regime shift. This provides a theoretical rationale, grounded in real
data, for the industry trend toward adaptive algorithmic smart order
routing.

Together, these results demonstrate that the multi-armed bandit framework
--- extended to the multi-agent setting with explicit collision modelling
--- is both structurally appropriate for the dark pool routing problem and
empirically informative when calibrated to real market data.

% ─────────────────────────────────────────────────────────────────────────────
\bibliographystyle{plainnat}
\bibliography{references}

\end{document}
